\documentclass[11pt]{article}
\usepackage[a4paper,margin=26mm]{geometry}
\usepackage[T1]{fontenc}
\usepackage{lmodern}
\usepackage{amsmath,amssymb,amsthm}
\usepackage{microtype}
\usepackage{array}
\usepackage{xurl}
\usepackage{hyperref}
\hypersetup{colorlinks=true,linkcolor=blue,urlcolor=blue,citecolor=blue,
 pdftitle={Medial racks and quandles of prime order},
 pdfauthor={Horace Dong}}
\newtheorem{theorem}{Theorem}
\newtheorem{lemma}{Lemma}
\newtheorem{proposition}{Proposition}
\newtheorem{corollary}{Corollary}
\theoremstyle{definition}
\newtheorem{definition}{Definition}
\theoremstyle{remark}
\newtheorem{remark}{Remark}
\newcommand{\repoversion}{1.25.0}
\setlength{\emergencystretch}{1.5em}
\newcommand{\F}{\mathbb F}
\newcommand{\rmed}{r_{\mathrm{med}}}
\newcommand{\rred}{r_{\mathrm{2\text{-}red}}}
\newcommand{\qmed}{q_{\mathrm{med}}}
\newcommand{\qred}{q_{\mathrm{2\text{-}red}}}
\DeclareMathOperator{\Dis}{Dis}
\DeclareMathOperator{\Aut}{Aut}
\DeclareMathOperator{\Sym}{Sym}

\title{Medial racks and quandles of prime order:\\ a proof of a conjecture of Vojt\v{e}chovsk\'y and Yang}
\author{Horace Dong\\ \small Independent researcher\\
\small\texttt{maxwelldhx@gmail.com}\\
\small ORCID: \href{https://orcid.org/0009-0002-7554-7548}{0009-0002-7554-7548}}
\date{22 November 2026}
\begin{document}
\maketitle

\begin{abstract}
A rack is a set with a binary operation whose left translations are
bijections and which satisfies the left self-distributive law; a quandle is an
idempotent rack. Vojt\v{e}chovsk\'y and Yang enumerated racks and quandles of
order at most~$13$ up to isomorphism and conjectured that, for every prime
$p$, the number of medial racks of order $p$ exceeds the number of
$2$-reductive racks of order $p$ by exactly $p-2$, that the same holds for
quandles, and that every medial rack of order $p$ that is not $2$-reductive is
a quandle. We prove this conjecture. More precisely, we show that every medial
rack of prime order $p$ that is not $2$-reductive is isomorphic to exactly one
of the affine quandles $x*y=(1-\lambda)x+\lambda y$ on $\F_p$ with
$\lambda\ne0,1$. Connectedness and idempotence are not assumed: they follow
from a divisibility argument for the displacement group, whose point
stabilizers all lie in one subgroup. The three statements of the conjecture
are formally verified in Lean~4. The proofs were found and formalized by AI
agents working under the author's direction; their role is described at the
end of the paper.
\end{abstract}

\section{Introduction}

A \emph{rack} is a nonempty set $X$ with a binary operation $*$ such that every
left translation
\[
L_x\colon X\to X,\qquad y\mapsto x*y ,
\]
is a bijection, and the left self-distributive law
\[
x*(y*z)=(x*y)*(x*z)
\]
holds for all $x,y,z\in X$. A rack is a \emph{quandle} if it is idempotent,
that is, $x*x=x$ for all $x$. A rack is \emph{medial} if
\[
(x*u)*(v*y)=(x*v)*(u*y)\qquad\text{for all }x,u,v,y\in X,
\]
and \emph{$2$-reductive} if
\begin{equation}\label{eq:2red}
(x*u)*v=(y*u)*v\qquad\text{for all }x,y,u,v\in X ;
\end{equation}
this is identity~(1.1) of Vojt\v{e}chovsk\'y and Yang~\cite[p.~2]{vy}. A rack
that is not $2$-reductive is called \emph{non-$2$-reductive}. Every
$2$-reductive rack is medial~\cite[p.~2]{vy} (see also Lemma~\ref{lem:redmed}).
An isomorphism between operations $*$ on $X$ and $*'$ on $X'$ is a bijection
$f\colon X\to X'$ with $f(x*y)=f(x)*'f(y)$. These are the definitions
of~\cite[Section~1]{vy}; no connectedness is assumed, and the right
translations need not be bijective.

For $n\ge1$ let $\rmed(n)$, $\rred(n)$, $\qmed(n)$ and $\qred(n)$ be the
numbers of medial racks, $2$-reductive racks, medial quandles and
$2$-reductive quandles of order $n$ up to isomorphism. Vojt\v{e}chovsk\'y and
Yang tabulated these numbers for $n\le13$~\cite[Tables~1 and~2]{vy} and, at the
end of their paper, stated:

\begin{quote}
\textbf{Conjecture 6.5 of~\cite{vy}.} Let $p$ be a prime. Then
$\rmed(p)-\rred(p)=p-2$ and $\qmed(p)-\qred(p)=p-2$. In particular, every
medial rack of order $p$ that is not $2$-reductive is a quandle.
\end{quote}

\noindent Their data for the primes up to $13$ are reproduced in
Table~\ref{tab:data}. We prove the conjecture.

\begin{table}[ht]
\centering\small
\begin{tabular}{@{}rrrrrrr@{}}
\hline
$p$ & $\rmed(p)$ & $\rred(p)$ & difference & $\qmed(p)$ & $\qred(p)$ & difference\\
\hline
2 & 2 & 2 & 0 & 1 & 1 & 0\\
3 & 6 & 5 & 1 & 3 & 2 & 1\\
5 & 68 & 65 & 3 & 18 & 15 & 3\\
7 & 1965 & 1960 & 5 & 251 & 246 & 5\\
11 & 35982366 & 35982357 & 9 & 1246488 & 1246479 & 9\\
13 & 25731050872 & 25731050861 & 11 & 466087635 & 466087624 & 11\\
\hline
\end{tabular}
\caption{The prime orders in Tables~1 and~2 of~\cite{vy}. These values are
not used in the proofs and are not part of the Lean files.}
\label{tab:data}
\end{table}

\begin{theorem}\label{thm:main}
For every prime $p$:
\begin{enumerate}
\item[(a)] $\rmed(p)=\rred(p)+(p-2)$;
\item[(b)] $\qmed(p)=\qred(p)+(p-2)$;
\item[(c)] every medial rack of order $p$ that is not $2$-reductive is a
quandle.
\end{enumerate}
\end{theorem}

\noindent Since every $2$-reductive rack is medial, $\rred(p)\le\rmed(p)$ and
$\qred(p)\le\qmed(p)$, so (a) and (b) are the two equations of the conjecture
written as additions. Theorem~\ref{thm:main} is a consequence of the following
classification. For a prime $p$ and $\lambda\in\F_p$ let $A_\lambda$ denote
$\F_p$ with the operation
\[
x*y=(1-\lambda)x+\lambda y .
\]

\begin{theorem}\label{thm:class}
Let $p$ be a prime. For every $\lambda\in\F_p\setminus\{0,1\}$, $A_\lambda$ is
a medial quandle that is not $2$-reductive, and $A_\lambda\cong A_\mu$ only if
$\lambda=\mu$. Every medial rack of order $p$ that is not $2$-reductive is
isomorphic to $A_\lambda$ for some $\lambda\in\F_p\setminus\{0,1\}$.
\end{theorem}

\noindent Theorem~\ref{thm:class} is proved in the Lean files (for operations
on $\{0,\dots,p-1\}$) but is not one of the frozen statements; see
Section~\ref{sec:lean}.

Thus the non-$2$-reductive medial racks of order $p$ fall into exactly $p-2$
isomorphism classes, all of them quandles; for $p=2$ there are none. The main
point is that the input is an arbitrary medial rack, possibly not connected and
not idempotent. The argument uses the displacement group $D$, which is abelian
for medial racks. Fix a left translation $T$. Every point stabilizer of $D$ is
contained in the centralizer $H$ of $T$ in $D$ (Lemma~\ref{lem:stab}), and
$H=D$ forces the rack to be $2$-reductive (Lemma~\ref{lem:top}). So for a
non-$2$-reductive medial rack the index $[D:H]>1$ divides the size of every
$D$-orbit. When the order is a prime $p$, the orbit sizes add up to $p$, so the
index is $p$ and $D$ is transitive, hence regular (Lemma~\ref{lem:prime}).
Coordinates from the cyclic group $D$ then put the operation in the form
$c+(1-\lambda)x+\lambda y$, and the rack axioms force $c=0$
(Proposition~\ref{prop:affine}).

\subsection*{Prior work}
The conjecture and the enumeration behind it are due to Vojt\v{e}chovsk\'y and
Yang~\cite{vy}, who relied in part on counts of $2$-reductive racks and
quandles due to Jedli\v{c}ka, Pilitowska, Stanovsk\'y and Zamojska-Dzienio:
the quandle counts are in~\cite{jpsz}, and the rack counts were supplied by
Jedli\v{c}ka using the methods of~\cite{jpsz}~\cite[p.~5]{vy}. That a rack is medial if and only if its
displacement group is abelian is~\cite[Proposition~2.4(ii)]{hsv}, as quoted
in~\cite[Section~1.1]{vy}; Lemma~\ref{lem:dis}(d) is the direction used here.
Etingof, Soloviev and Guralnick~\cite{esg} showed that the indecomposable
non-degenerate set-theoretical solutions of the quantum Yang--Baxter equation
on a set of prime order are affine; in the language of quandles, every
connected quandle of prime order is affine (a quandle is connected if the group
generated by its left translations acts transitively~\cite[p.~3]{vy}). Hulpke, Stanovsk\'y and
Vojt\v{e}chovsk\'y state this as~\cite[Theorem~9.3]{hsv}, attribute it
to~\cite{esg} and give two new proofs. They also note that two connected affine
quandles over the same abelian group $A$ are isomorphic if and only if the
defining automorphisms are conjugate in $\Aut(A)$~\cite[p.~14]{hsv}; for
$A=\F_p$ this criterion, translated from right to left translations, gives the
uniqueness part of Theorem~\ref{thm:class}, which we prove directly in
Lemma~\ref{lem:models}(b). By
Theorem~\ref{thm:class} and Corollary~\ref{cor:regular}, the classes counted by
the two differences are represented by connected affine quandles, which fall
under the classification of~\cite{esg,hsv}. What the conjecture needs in
addition is that a medial rack of prime order which is not $2$-reductive is
connected and idempotent; the divisibility argument above supplies this. For
medial quandles, Jedli\v{c}ka et al.\ prove conditions on the orbit sizes that
force $m$-reductivity~\cite[Corollaries~6.3 and~6.10]{jpsz}; these concern
quandles, so they do not apply to the racks counted in (a); and even for
quandles of prime order they do not give $2$-reductivity directly, since
Corollary~6.3 yields only $3$-reductivity and the hypothesis of
Corollary~6.10 can fail (orbit sizes $6$, $10$, $15$ sum to the prime $31$). As far as we know,
Conjecture~6.5 of~\cite{vy} has not been proved elsewhere; the searches are
described in the final section.

\section{Preliminaries}\label{sec:prelim}

Throughout, permutations are composed as functions: $(gh)(z)=g(h(z))$.
Vojt\v{e}chovsk\'y and Yang write permutations on the right,
$yL_x=x*y$, and define the displacement group as
$\Dis(X)=\langle L_x^{-1}L_y : x,y\in X\rangle$~\cite[Section~1.1]{vy}; in
function notation the generator $L_x^{-1}L_y$ (first $L_x^{-1}$, then $L_y$) is
$L_yL_x^{-1}$. Hulpke, Stanovsk\'y and Vojt\v{e}chovsk\'y work with right
translations instead~\cite{hsv}. The Lean statements are identities in the
operation itself, so none of these conventions enters them.

Let $(X,*)$ be a rack. The self-distributive law says that every $L_x$ is an
automorphism of $(X,*)$. If $g$ is an automorphism and $x\in X$, then
\begin{equation}\label{eq:conj}
gL_xg^{-1}=L_{g(x)},
\end{equation}
since $g(x*g^{-1}(y))=g(x)*y$. In particular $L_{x*u}=L_xL_uL_x^{-1}$.

\begin{lemma}\label{lem:redmed}
Every $2$-reductive rack is medial.
\end{lemma}

\begin{proof}
By~\eqref{eq:2red} with $y=u$, and by $L_{u*u}=L_uL_uL_u^{-1}=L_u$, we get
$(x*u)*v=(u*u)*v=u*v$ for all $x,u,v$. Hence
\[
(x*u)*(v*y)=u*(v*y)=(u*v)*(u*y)=v*(u*y)=(x*v)*(u*y).\qedhere
\]
\end{proof}

\begin{lemma}\label{lem:inv}
Being a rack, a quandle, medial or $2$-reductive is preserved by isomorphisms.
\end{lemma}

\begin{proof}
Each property is a set of identities in the operation, together with the
bijectivity of the left translations; an isomorphism $f$ transports both,
since $L'_{f(x)}=fL_xf^{-1}$.
\end{proof}

\section{The displacement group of a medial rack}\label{sec:dis}

Let $(X,*)$ be a rack, fix a base point $b\in X$, write $T=L_b$, and let
\[
D=\bigl\langle L_uT^{-1} : u\in X\bigr\rangle\le\Sym(X).
\]
Since $L_uL_v^{-1}=(L_uT^{-1})(L_vT^{-1})^{-1}$, the group $D$ contains every
$L_uL_v^{-1}$, so $D=\Dis(X)$ in the notation of~\cite{vy}.

\begin{lemma}\label{lem:dis}
\begin{enumerate}
\item[(a)] Every element of $D$ is an automorphism of $(X,*)$.
\item[(b)] $L_xdL_x^{-1}\in D$ for every $x\in X$ and every $d\in D$.
\item[(c)] $L_uL_v^{-1}\in D$ for all $u,v\in X$.
\item[(d)] If $(X,*)$ is medial, then $D$ is abelian.
\end{enumerate}
\end{lemma}

\begin{proof}
(a) The generators are products of left translations and their inverses.
(c) was shown above. (b) By~\eqref{eq:conj},
$L_x(L_uT^{-1})L_x^{-1}=L_{x*u}L_{x*b}^{-1}$, which lies in $D$ by (c). Since
conjugation by $L_x$ is a group homomorphism, it maps the subgroup generated by
the elements $L_uT^{-1}$ into $D$. (d) The medial law says $L_{x*u}L_v=L_{x*v}L_u$ (apply both sides to $y$).
Take $x=b$ and substitute $L_{b*u}=TL_uT^{-1}$ and $L_{b*v}=TL_vT^{-1}$
from~\eqref{eq:conj}: $TL_uT^{-1}L_v=TL_vT^{-1}L_u$. Multiplying by $T^{-1}$ on
the left and on the right gives
$(L_uT^{-1})(L_vT^{-1})=(L_vT^{-1})(L_uT^{-1})$. So the generators of $D$
commute, and $D$ is abelian.
\end{proof}

Part (d) is the direction ``medial $\Rightarrow$ $\Dis$ abelian''
of~\cite[Proposition~2.4(ii)]{hsv}.

From now on let $(X,*)$ be a medial rack, and let
\[
H=\{d\in D : dT=Td\}
\]
be the centralizer of $T$ in $D$. (Since $D$ is abelian and
$L_x=(L_xT^{-1})T$, conjugation by any $L_x$ induces the same automorphism
$f$ of $D$, and $H$ is its fixed-point subgroup; this description is not
needed below.) For $x\in X$ let $D_x=\{d\in D: d(x)=x\}$ be the stabilizer of
$x$ in $D$.

\begin{lemma}\label{lem:stab}
$D_x\le H$ for every $x\in X$.
\end{lemma}

\begin{proof}
Let $d\in D_x$. By Lemma~\ref{lem:dis}(a) and~\eqref{eq:conj},
$dL_xd^{-1}=L_{d(x)}=L_x$, so $d$ commutes with $L_x$. By
Lemma~\ref{lem:dis}(d), $d$ commutes with $L_xT^{-1}\in D$. Hence $d$ commutes
with $(L_xT^{-1})^{-1}L_x=T$.
\end{proof}

\begin{lemma}\label{lem:top}
If $H=D$, then $(X,*)$ is $2$-reductive.
\end{lemma}

\begin{proof}
If $T$ commutes with every element of $D$, then the left translations
$L_x=(L_xT^{-1})T$ commute pairwise, because $D$ is abelian. Then
$L_{x*u}=L_xL_uL_x^{-1}=L_u$ for all $x,u$, so
$(x*u)*v=u*v=(y*u)*v$.
\end{proof}

\section{Prime order: the displacement group is regular}\label{sec:prime}

\begin{lemma}\label{lem:prime}
Let $G$ be a group of permutations of a set $X$ with $p$ elements, $p$ prime,
and let $K$ be a proper subgroup of $G$ that contains the stabilizer $G_x$ of
every point $x\in X$. Then $G$ is transitive on $X$. If moreover $G$ is
abelian, then for every $x\in X$ the map $G\to X$, $g\mapsto g(x)$, is a
bijection; in particular $|G|=p$.
\end{lemma}

\begin{proof}
Since $G_x\le K$, the index $[G:K]$ divides $[G:G_x]$, which is the size of
the orbit of $x$. The orbits partition $X$, so $[G:K]$ divides their total
size $p$. As $K\ne G$, we get $[G:K]=p$. Hence every orbit has at least $p$
elements, and there is only one orbit. Now let $G$ be abelian, and let
$g,h\in G$ with $g(x)=h(x)$. Every $y\in X$ has the form $y=k(x)$ with $k\in
G$, and $g(y)=k(g(x))=k(h(x))=h(y)$. So $g=h$; the map is injective, and it is
surjective by transitivity.
\end{proof}

\begin{corollary}\label{cor:regular}
Let $(X,*)$ be a medial rack of prime order $p$ that is not $2$-reductive.
Then $D$ is abelian, $|D|=p$, and for every $x\in X$ the map $d\mapsto d(x)$
is a bijection $D\to X$.
\end{corollary}

\begin{proof}
By Lemma~\ref{lem:top}, $H$ is a proper subgroup of $D$; by
Lemma~\ref{lem:stab} it contains every point stabilizer. Apply
Lemma~\ref{lem:prime} with $G=D$ and $K=H$, using Lemma~\ref{lem:dis}(d).
\end{proof}

\noindent Nowhere in this section or the previous one is the rack assumed to be
connected or idempotent; transitivity of $D$ is a conclusion.

\section{Affine coordinates}\label{sec:affine}

\begin{proposition}\label{prop:affine}
Let $(X,*)$ be a medial rack of prime order $p$ that is not $2$-reductive.
Then $(X,*)\cong A_\lambda$ for some $\lambda\in\F_p\setminus\{0,1\}$.
\end{proposition}

\begin{proof}
By Corollary~\ref{cor:regular}, $D$ has order $p$ and is therefore cyclic. Fix
an isomorphism $t\mapsto d_t$ from $(\F_p,+)$ onto $D$, so that
$d_{s+t}=d_sd_t$, and fix a point $o\in X$. By Corollary~\ref{cor:regular} the
map $e\colon\F_p\to X$, $e(t)=d_t(o)$, is a bijection, and
$d_t(e(y))=e(t+y)$. Transport the operation: for $x,y\in\F_p$ put
$x\circ y=e^{-1}\bigl(e(x)*e(y)\bigr)$. Then $e$ is an isomorphism from
$(\F_p,\circ)$ onto $(X,*)$, so $(\F_p,\circ)$ is a non-$2$-reductive rack
(Lemma~\ref{lem:inv}). It has two further properties.

\emph{Translation invariance.} Since $d_t$ is an automorphism of $(X,*)$,
\[
(t+x)\circ(t+y)=t+(x\circ y)\qquad\text{for all }t,x,y.
\]

\emph{Normalization.} Let $L=L_{e(0)}$. By Lemma~\ref{lem:dis}(b),
$Ld_tL^{-1}\in D$, say $Ld_tL^{-1}=d_{w}$ with $w=w(t)\in\F_p$. Applying both sides to $L(e(y))$ gives
$L(e(t+y))=e(w+(0\circ y))$, that is,
\[
0\circ(t+y)=w+(0\circ y)\qquad\text{for all }y .
\]

Let $c=0\circ0$ and $\phi(t)=0\circ t-c$. Taking $y=0$ above gives
$0\circ t=w+c$, hence $\phi(t+y)=w+0\circ y-c=\phi(t)+\phi(y)$; so $\phi$ is
additive on $\F_p$ and $\phi(t)=\lambda t$ with $\lambda=\phi(1)$. Thus
$0\circ y=\lambda y+c$, and by translation invariance
\[
x\circ y=x+\bigl(0\circ(y-x)\bigr)=c+(1-\lambda)x+\lambda y .
\]
If $\lambda=1$, then $x\circ y=y+c$ does not depend on $x$, and
$(x\circ u)\circ v=v+c=(y\circ u)\circ v$, contrary to non-$2$-reductivity; so
$\lambda\ne1$. If $\lambda=0$, then $0\circ0=0\circ1$, contrary to the
injectivity of left translations; so $\lambda\ne0$. Finally, the
self-distributive law with $x=y=z=0$ gives $0\circ c=c\circ c$, that is,
$\lambda c+c=2c$, so $(1-\lambda)c=0$ and $c=0$. Hence $(\F_p,\circ)=A_\lambda$
and $(X,*)\cong A_\lambda$.
\end{proof}

\begin{lemma}\label{lem:models}
Let $p$ be a prime and $\lambda\in\F_p\setminus\{0,1\}$.
\begin{enumerate}
\item[(a)] $A_\lambda$ is a rack, it is idempotent and medial, and it is not
$2$-reductive.
\item[(b)] If $\mu\in\F_p$ and $A_\lambda\cong A_\mu$, then $\lambda=\mu$.
\end{enumerate}
\end{lemma}

\begin{proof}
(a) $L_x(y)=(1-\lambda)x+\lambda y$ is a bijection because $\lambda\ne0$. Both
sides of the self-distributive law equal
$(1-\lambda)x+\lambda(1-\lambda)y+\lambda^2z$; both sides of the medial law
equal $(1-\lambda)^2x+\lambda(1-\lambda)(u+v)+\lambda^2y$; and
$x*x=x$. For $x=1$ and $y=u=v=0$, the two sides of~\eqref{eq:2red} are
$(1-\lambda)^2$ and $0$, which differ because $\lambda\ne1$.

(b) Let $f\colon\F_p\to\F_p$ be a bijection with
$f\bigl((1-\lambda)x+\lambda y\bigr)=(1-\mu)f(x)+\mu f(y)$. Putting $y=0$,
respectively $x=0$,
\[
f\bigl((1-\lambda)x\bigr)=(1-\mu)f(x)+\mu f(0),\qquad
f(\lambda y)=(1-\mu)f(0)+\mu f(y).
\]
Let $g(x)=f(x)-f(0)$. Given $a,b\in\F_p$, write $a=(1-\lambda)x$ and
$b=\lambda y$, which is possible because $1-\lambda$ and $\lambda$ are
nonzero; the three displayed identities give $g(a+b)=g(a)+g(b)$. So $g$ is
additive and $g(x)=x\,g(1)$, where $g(1)\ne0$ because $f$ is injective. The
second identity with $y=1$ gives $g(\lambda)=\mu\,g(1)$, while
$g(\lambda)=\lambda\,g(1)$. Hence $\lambda=\mu$.
\end{proof}

\begin{proof}[Proof of Theorem~\ref{thm:class}]
Combine Lemma~\ref{lem:models} and Proposition~\ref{prop:affine}.
\end{proof}

\section{The two counts}\label{sec:count}

\begin{proof}[Proof of Theorem~\ref{thm:main}]
Part (c) follows from Theorem~\ref{thm:class}: a non-$2$-reductive medial rack
of order $p$ is isomorphic to some $A_\lambda$, which is idempotent, and
idempotence is preserved by isomorphisms.

For (a), every rack of order $p$ is isomorphic to a rack on a fixed set of $p$
elements, so $\rmed(p)$ is the number of isomorphism classes of medial racks on
that set, and similarly for the other counts. By Lemmas~\ref{lem:redmed}
and~\ref{lem:inv}, the isomorphism classes of medial racks split into the
classes of $2$-reductive racks and the classes of non-$2$-reductive medial
racks. By Theorem~\ref{thm:class}, $\lambda\mapsto[A_\lambda]$ is a bijection
from $\F_p\setminus\{0,1\}$ onto the latter: it lands there by
Lemma~\ref{lem:models}(a), it is injective by Lemma~\ref{lem:models}(b), and
it is surjective by Proposition~\ref{prop:affine}. Hence
$\rmed(p)=\rred(p)+(p-2)$.

For (b), the same splitting holds among quandles, since a $2$-reductive
quandle is a medial quandle. Every $A_\lambda$ is idempotent, so all $p-2$
parameters already give quandles, and no parameter is lost when one restricts
to idempotent operations. Conversely, a non-$2$-reductive medial quandle is in
particular a non-$2$-reductive medial rack and is covered by
Proposition~\ref{prop:affine}. The isomorphism condition does not depend on
whether one regards the operations as racks or as quandles. So
$\lambda\mapsto[A_\lambda]$ is again a bijection onto the classes of
non-$2$-reductive medial quandles, and $\qmed(p)=\qred(p)+(p-2)$. For $p=2$
the parameter set is empty and both differences are~$0$.
\end{proof}

{\sloppy
\section{Formal verification}\label{sec:lean}

Theorem~\ref{thm:main} is formally verified in Lean~4 (version 4.34.1) with
Mathlib~\cite{mathlib} (commit \texttt{d13f23b7}). Theorem~\ref{thm:class} is
proved in the same files, for operations on $\{0,\dots,p-1\}$, on the way to
Theorem~\ref{thm:main}, but it is not among the frozen statements. The $13$
files are in version~\repoversion{} of the repository~\cite{repo}, directory
\texttt{lean/Research/TaskP/PrimeMedialRacks/}.

\paragraph{The statements.}
The statement file \texttt{Challenge.lean} contains the definitions and the
three statements; its only proof is the short verification, inside
\texttt{isoSetoid}, that isomorphism is an equivalence relation, which is
needed to form the quotient; \texttt{classCount} is defined by a three-line
tactic block. In ASCII transliteration (the module comment, the
docstrings and the \texttt{import}, \texttt{namespace} and \texttt{end} lines
omitted; the file itself uses Unicode):
{\footnotesize
\begin{verbatim}
abbrev Operation (alpha : Type*) := alpha -> alpha -> alpha
def IsRack {alpha : Type*} (op : Operation alpha) : Prop :=
  (forall x, Function.Bijective (op x)) /\
    forall x y z, op x (op y z) = op (op x y) (op x z)
def Idempotent {alpha : Type*} (op : Operation alpha) : Prop :=
  forall x, op x x = x
def Medial {alpha : Type*} (op : Operation alpha) : Prop :=
  forall x u v y, op (op x u) (op v y) = op (op x v) (op u y)
def TwoReductive {alpha : Type*} (op : Operation alpha) : Prop :=
  forall x y u v, op (op x u) v = op (op y u) v
def Isomorphic {alpha beta : Type*} (op : Operation alpha) (op' : Operation beta) : Prop :=
  exists e : alpha ~= beta, forall x y, e (op x y) = op' (e x) (e y)
def isoSetoid (n : Nat) : Setoid (Operation (Fin n)) where
  r := Isomorphic
  iseqv := {
    refl := fun _ => <Equiv.refl _, fun _ _ => rfl>
    symm := by
      rintro _ _ <e, h>
      refine <e.symm, fun x y => e.injective ?_>
      simp [h]
    trans := by
      rintro _ _ _ <e, h> <f, k>
      exact <e.trans f, fun x y => by simp [h, k]> }
abbrev IsoClass (n : Nat) := Quotient (isoSetoid n)
noncomputable def classCount (n : Nat) (P : Operation (Fin n) -> Prop) : Nat := by
  classical
  let C := {q : IsoClass n // exists op, Quotient.mk (isoSetoid n) op = q /\ P op}
  letI : Fintype C := Fintype.ofFinite C
  exact Fintype.card C
def RackCount : Prop :=
  forall p : Nat, p.Prime ->
    classCount p (fun op => IsRack op /\ Medial op) =
      classCount p (fun op => IsRack op /\ TwoReductive op) + (p - 2)
def QuandleCount : Prop :=
  forall p : Nat, p.Prime ->
    classCount p (fun op => IsRack op /\ Idempotent op /\ Medial op) =
      classCount p (fun op => IsRack op /\ Idempotent op /\ TwoReductive op) + (p - 2)
def AutomaticIdempotence : Prop :=
  forall (p : Nat), p.Prime -> forall op : Operation (Fin p),
    IsRack op -> Medial op -> not TwoReductive op -> Idempotent op
\end{verbatim}}
\noindent Here \texttt{alpha}, \texttt{beta} stand for $\alpha$, $\beta$,
\texttt{forall}, \texttt{exists}, \texttt{/\textbackslash}, \texttt{->} and
\texttt{not} for $\forall$, $\exists$, $\wedge$, $\to$, $\neg$,
\texttt{\textasciitilde=} for $\simeq$ (a bijection, Lean's \texttt{Equiv}),
\texttt{<}\dots\texttt{>} for the anonymous constructor
$\langle\dots\rangle$, \texttt{Nat} for $\mathbb N$, \texttt{Fin n} is the set
$\{0,\dots,n-1\}$, and \texttt{p - 2} is subtraction of natural numbers.
\texttt{op x} is the left translation $L_x$, so \texttt{IsRack},
\texttt{Idempotent}, \texttt{Medial} and \texttt{TwoReductive} are the
definitions of Section~1 with the variables in the order of~\cite{vy}.
\texttt{Isomorphic} relabels both inputs and the output by one bijection.
\texttt{classCount n P} is the number of isomorphism classes of
\emph{all} binary operations on \texttt{Fin n} that contain an operation
satisfying \texttt{P}. All four properties are invariant under isomorphism
(Lemma~\ref{lem:inv}; \texttt{iso\_\allowbreak rack} and its companions in
\texttt{IsoTransfer.lean}), so ``some member satisfies \texttt{P}'' is the
same as ``every member does'', and since every rack of order $p$ is isomorphic
to an operation on \texttt{Fin p}, \texttt{classCount p (fun op => IsRack op
/\textbackslash{} Medial op)} is $\rmed(p)$; likewise for the other three
counts. \texttt{RackCount} and \texttt{QuandleCount} are parts (a) and (b) of
Theorem~\ref{thm:main}, that is, the two equations of the conjecture in
additive form; as noted after Theorem~\ref{thm:main}, the two forms are
equivalent, and for a prime $p$ the subtraction \texttt{p - 2} is the ordinary
one. Both statements cover every prime, including $p=2$.
\texttt{AutomaticIdempotence} is part (c), the sentence that follows the two
equations in~\cite{vy}, stated for all operations on \texttt{Fin p}. None of
the statements assumes connectedness, idempotence of the racks, or an affine
form. The final reviewer (see the last section) checked these points against
the source and found no weakening, extra hypothesis or vacuity.

\paragraph{Review and freezing.}
\texttt{Challenge.lean} was reviewed against~\cite{vy} by a Codex agent that
had written neither the proofs nor the Lean statement, and frozen, by recording
its SHA-256 hash, on 20 November 2026 at 18:04~UTC. The development logs record
no compilation of a proof file before that time, and the file was not changed
afterwards.

\paragraph{The proofs.}
In \texttt{Proof.lean}, \texttt{Proof.classified} proves that every
non-$2$-reductive medial rack on \texttt{Fin p} is isomorphic to some
$A_\lambda$ (the property \texttt{ClassifiedAt p} of
\texttt{Enumeration.lean}); \texttt{Proof.RackCount},
\texttt{Proof.QuandleCount} and \texttt{Proof.\allowbreak AutomaticIdempotence}
apply \texttt{rack\_\allowbreak count\_\allowbreak of\_\allowbreak classification},
\texttt{quandle\_\allowbreak count\_\allowbreak of\_\allowbreak classification}
and \texttt{idempotent\_\allowbreak of\_\allowbreak classification} to it.
Table~\ref{tab:lean} lists the Lean counterparts of the steps of the written
proof (names in the namespace \texttt{TaskP.PrimeMedialRacks}). In Lean the
displacement group is the subgroup \texttt{displacement} generated by the
permutations $L_uL_b^{-1}$, as in Section~\ref{sec:dis}, and $\lambda$ is
written \texttt{k}.

\begin{table}[ht]
\centering\footnotesize
\renewcommand{\arraystretch}{1.05}
\begin{tabular}{@{}>{\raggedright\arraybackslash}p{0.44\textwidth}>{\raggedright\arraybackslash}p{0.52\textwidth}@{}}
\hline
Step of the written proof & Lean (file) \\
\hline
Lemma~\ref{lem:redmed} & \texttt{reductive\_medial} (\texttt{IsoTransfer}) \\
Lemma~\ref{lem:inv} & \texttt{iso\_rack}, \texttt{iso\_idempotent},
\texttt{iso\_medial}, \texttt{iso\_reductive}, \texttt{iso\_symm},
\texttt{iso\_trans} (\texttt{IsoTransfer}) \\
Left translations as permutations; identity~\eqref{eq:conj} &
\texttt{leftPerm} (\texttt{DisplacementData}); \texttt{leftPerm\_mem\_aut},
\texttt{conjugate\_leftPerm} (\texttt{DisplacementGroup}) \\
Lemma~\ref{lem:dis}: the group $D$; (a), (b), (c), (d) &
\texttt{displacement}, \texttt{displacement\_acts},
\texttt{displacement\_normalizes}, \texttt{difference\_mem},
\texttt{displacement\_commutes} (\texttt{DisplacementGroup}) \\
The subgroup $H$; Lemmas~\ref{lem:stab} and~\ref{lem:top} &
\texttt{fixedDisplacements}, \texttt{stabilizer\_le\_fixed},
\texttt{fixed\_ne\_top} (\texttt{PrimeOrbits}) \\
Lemma~\ref{lem:prime} & \texttt{prime\_action\_transitive},
\texttt{abelian\_transitive\_regular} (\texttt{PrimeAction}) \\
Corollary~\ref{cor:regular} & \texttt{nonempty\_regularDisplacements}
(\texttt{PrimeOrbits}), with the structure \texttt{RegularDisplacements}
(\texttt{DisplacementData}) \\
Proposition~\ref{prop:affine}: coordinates from the cyclic group $D$ &
\texttt{regular\_affine} (\texttt{Coordinates}) \\
Proposition~\ref{prop:affine}: the affine form, $\lambda\ne0,1$, $c=0$ &
\texttt{normalized\_coordinates} (\texttt{CoordinatesCore}) \\
The operation of $A_\lambda$; Lemma~\ref{lem:models}(a) & \texttt{affine},
\texttt{affine\_isRack}, \texttt{affine\_idempotent},
\texttt{affine\_medial}, \texttt{affine\_nonreductive} (\texttt{Affine}) \\
Lemma~\ref{lem:models}(b) & \texttt{affine\_parameter\_unique}
(\texttt{Affine}) \\
$A_\lambda$ moved to $\{0,\dots,p-1\}$ & \texttt{labels},
\texttt{affineTable}, \texttt{affineTable\_iso} (\texttt{ClassCount});
\texttt{affineTable\_properties} (\texttt{Enumeration}) \\
Splitting the classes; $|\F_p\setminus\{0,1\}|=p-2$; the bijection with the
parameters & \texttt{classCount\_partition}, \texttt{parameter\_card},
\texttt{count\_affine\_models} (\texttt{ClassCount}) \\
Theorem~\ref{thm:main}(a), (b), (c) from the classification &
\texttt{rack\_count\_of\_classification},
\texttt{quandle\_count\_of\_classification},
\texttt{idempotent\_of\_classification} (\texttt{Enumeration}) \\
\hline
\end{tabular}
\caption{Lean counterparts of the steps of the proof.}
\label{tab:lean}
\end{table}

\noindent The finite group theory comes from Mathlib: the decomposition of a
finite set into orbits, the orbit--stabilizer correspondence, divisibility of
indices, and the fact that a group of prime order is cyclic. The test file
\texttt{CheckRack.lean} restates the three theorems with the frozen types
(\texttt{check\_\allowbreak RackCount}, \texttt{check\_\allowbreak
QuandleCount}, \texttt{check\_\allowbreak AutomaticIdempotence}). The axioms of
these three theorems and of \texttt{Proof.RackCount},
\texttt{Proof.QuandleCount} and \texttt{Proof.AutomaticIdempotence}, printed by
\texttt{CheckRack.lean}, are exactly \texttt{propext},
\texttt{Classical.choice} and \texttt{Quot.sound}. No file uses
\texttt{sorry}, \texttt{decide}, \texttt{native\_\allowbreak decide},
\texttt{bv\_\allowbreak decide} or an \texttt{axiom} declaration.

\paragraph{Checks.}
The $13$ files ($10$ proof modules, \texttt{Challenge.lean},
\texttt{Proof.lean} and \texttt{CheckRack.lean}) were compiled module by
module in two new empty directories on the author's Mac, with the pinned
toolchain and with warnings treated as errors, on 20 November 2026
(18:16--18:21~UTC), and every module was replayed with \texttt{leanchecker};
each directory also ran one build of the whole package. All $54$ commands
($27$ per directory) ended with exit code~$0$, the \texttt{leanchecker} runs
printed nothing, and the two directories produced identical \texttt{.olean}
files. The final reviewer (see the next section) compiled the $13$ files again,
together with a short audit file of its own, in a third new directory on
20 November 2026, calling \texttt{lean} directly with warnings treated as
errors, and replayed every module with \texttt{leanchecker}; all runs ended
with exit code~$0$ (according to its report, between 18:32 and 18:35~UTC, while
it held the machine's Lean lock). Its audit file checks the three frozen
statements against the test theorems with fully qualified names. Mathlib and
its dependencies were taken from a shared prebuilt cache at the pinned commits
and were not rechecked; the \texttt{leanchecker} replays cover the $13$ project
modules. In the repository the files have module names under
\texttt{Research.TaskP.PrimeMedialRacks} and differ from the checked files only
in their \texttt{import} lines (\texttt{Challenge.lean} is unchanged).
The repository versions were made by a script that changed only the local \texttt{import} lines, and on 20 November 2026 (19:55--20:07~UTC) they were compiled and replayed with \texttt{leanchecker} again, module by module, in two new empty directories on the author's Mac, one after the other, together with the files of three other papers of version~\repoversion{} ($39$ modules in all); there were no warnings, and the two directories produced identical \texttt{.olean} files and the same axiom lists. In the whole-project rebuild of version~\repoversion{} of the repository (on the author's Windows computer, in two passes on 20 and 21 November 2026), these files were compiled and replayed with \texttt{leanchecker} again in pass~2 (21 November 2026), which built the files not included in pass~1 on top of its results; no module needed a second attempt (details in the repository). The formalization has not been checked with
comparator or with an independent kernel.

\par}

\section*{How this paper was produced}
This work was carried out from 20 to 22 November 2026 (UTC) in a workflow directed by
the author, using OpenAI Codex with the model GPT-6.2-Astra (reasoning level
ultra) and Anthropic Claude Code with the model Claude Opus~5.5. The author
asked for a large new search for open problems, and Claude Code passed this
request to Codex at 17:23~UTC. Between 17:23 and 17:53~UTC Codex screened
$111$ problems from the literature and wrote complete proofs, each read by a
second agent, for the five top-ranked candidates; the
proof for this problem (the divisibility argument and the affine
classification) was written by 17:29~UTC, according to the file times, and was
read by another Codex agent (same model) that had not taken part in it. While
the search was still running, the author asked that the five most important
candidates be formalized at once, without waiting for the author's choice;
Claude Code passed this request to Codex at 17:48~UTC. At 18:00~UTC Codex
ranked the candidates by importance and put this problem first. From then on,
Codex agents (same model) checked the reading of the source, searched again for
prior work, wrote the statement file, the Lean proof and the written proof,
and ran the checks. The reading of the source was checked by a Codex agent
that had taken part neither in the proof nor in drafting the Lean statement,
and the statement file was reviewed before it was frozen by a Codex agent that
had written neither the proofs nor the Lean statement. According to the file
timestamps and the development logs, the proof files were created after the
freeze, between 18:05 and 18:15~UTC, and compiled between 18:05 and
18:16~UTC; they follow the proof written in the search for problems. After the
checks, an independent Claude Code agent (model Claude Opus~5.5) that had
written none of the files gave the final review. It read the source itself,
compared the statement file with it, read the written proof, the proof files
and the raw logs of the checks, recomputed by brute force in Python all racks
and quandles of order at most~$5$ up to isomorphism (reproducing the
corresponding entries of~\cite[Tables~1 and~2]{vy}), and rebuilt the
formalization as described above. Its decision was APPROVE\_FORMAL, with five
minor remarks on the written account (one optional), all taken into account
here. Claude Code (Claude Opus~5.5) checked the sources and wrote this text
from Codex's manuscript and the Lean files. An independent Claude Code agent (model Claude Opus~5.5) that had written none of the files reviewed this paper, and another reviewed the release; their corrections were applied. Throughout, ``we''
refers to this workflow; the reading and searching described in the paper were
done by the agents. Codex agents searched for prior work on 20 November 2026,
in the search for problems and again before the statement was frozen; they
found no proof of Conjecture~6.5 of~\cite{vy} (details in the repository).
Claude Code searched again on 21 November 2026 (17:54--18:21~UTC; the arXiv, Semantic Scholar, Crossref, DataCite, Zenodo, MathDB, GitHub and repositories of AI-assisted mathematics) and found no new related work; that day the citation index OpenAlex refused to list the works citing~\cite{vy} (rate limit), and the works citing it that two other indexes list concern other questions. Claude Code searched again on 22 November 2026 (05:39--05:51~UTC), when the OpenAlex lists of citing works refused on 21 November were retrieved, and found nothing new. This reports the coverage of the searches, not a guarantee
that they were complete. The author is not a professional mathematician.
The author has read this paper and, using a side-by-side table, compared its Lean statements with the main theorem, Theorem~\ref{thm:main}. No human mathematician has reviewed this paper. The correctness
of the results rests on the Lean proofs and on the agreement of the Lean
statements with the conjecture of~\cite{vy}; the written proofs are given for
the reader. AI systems are not authors of this paper; the author takes full
responsibility for its content.

\medskip
\noindent Thanks are due to Petr Vojt\v{e}chovsk\'y and Seung Yeop Yang for the
conjecture.

{\small
\begin{thebibliography}{9}
\setlength{\itemsep}{1pt}
\bibitem{repo} H.~Dong, \emph{ai4math-results}, repository,
\url{https://github.com/Horace-Maxwell/ai4math-results}, version~\repoversion{}
(2026); archived at Zenodo under the concept DOI
\url{https://doi.org/10.5281/zenodo.22970254}.
\bibitem{esg} P.~Etingof, A.~Soloviev and R.~Guralnick, Indecomposable
set-theoretical solutions to the quantum Yang--Baxter equation on a set with a
prime number of elements, \emph{J. Algebra} 242 (2001), no.~2, 709--719,
\url{https://doi.org/10.1006/jabr.2001.8842}; arXiv:math/0007170.
\bibitem{hsv} A.~Hulpke, D.~Stanovsk\'y and P.~Vojt\v{e}chovsk\'y, Connected
quandles and transitive groups, \emph{J. Pure Appl. Algebra} 220 (2016),
no.~2, 735--758, \url{https://doi.org/10.1016/j.jpaa.2015.07.014};
arXiv:1409.2249v2 (2015). Page numbers refer to the arXiv version.
\bibitem{jpsz} P.~Jedli\v{c}ka, A.~Pilitowska, D.~Stanovsk\'y and
A.~Zamojska-Dzienio, The structure of medial quandles, \emph{J. Algebra} 443
(2015), 300--334, \url{https://doi.org/10.1016/j.jalgebra.2015.04.046};
arXiv:1409.8396v2 (2015).
\bibitem{mathlib}
The mathlib Community, The Lean mathematical library, in \emph{Proc.\ 9th ACM
SIGPLAN Int.\ Conf.\ Certified Programs and Proofs} (CPP '20), ACM, 2020,
\url{https://doi.org/10.1145/3372885.3373824}.
\bibitem{vy} P.~Vojt\v{e}chovsk\'y and S.~Y. Yang, Enumeration of racks and
quandles up to isomorphism, \emph{Math. Comp.} 88 (2019), no.~319,
2523--2540, \url{https://doi.org/10.1090/mcom/3409}; arXiv:1911.04991v1
(2019). Page numbers refer to the arXiv version.
\end{thebibliography}}
\end{document}
